\subsection{Categorias e casos negativos}
\label{sec:categorias_testes}

O Quadro~\ref{quad:categorias_testes} organiza a suíte por responsabilidade.
As categorias se sobrepõem, pois um teste FFI pode exercitar várias fases do
motor e um teste de widget pode verificar uma regra do estado. Por isso, elas
não particionam a contagem total de testes.

\begin{quadro}[htb]
\centering
\legenda[quad:categorias_testes]{Categorias de testes e exemplos de condições negativas}
\footnotesize
\begin{tabular}{|P{2.1cm}|P{5.2cm}|P{5.5cm}|}
\hline
Categoria & Responsabilidade e arquivo de referência & Exemplo negativo verificado \\
\hline\hline
Lexer & Tokens e posições, \path{rust_core/src/lexer/mod.rs}. & Caractere desconhecido produz diagnóstico. \\
\hline
Parser & AST e recuperação, \path{rust_core/src/parser/mod.rs}. & Programa malformado retorna erros; tentativa de continuar após falha. \\
\hline
Semântico & Nomes e regras estáticas, \path{rust_core/src/semantic/mod.rs}. & Nome não declarado, aridade incorreta, controle fora do contexto e divisor constante zero. \\
\hline
Interpretador & Valores, escopos e rastro, \path{rust_core/src/interpreter/}. & Divisão por zero em execução e laço interrompido pelo limite configurado. \\
\hline
API & Contrato integrado, \path{rust_core/src/api.rs}. & Entrada sintaticamente inválida produz resultado sem sucesso. \\
\hline
Domínio & Documentos e progresso, \path{test/unit/workspace_test.dart} e \path{test/unit/learning_progress_test.dart}. & Entradas corrompidas, identificadores desconhecidos e lacunas na progressão. \\
\hline
Repositórios & Persistência, \path{test/unit/settings_test.dart} e \path{test/unit/learning_progress_test.dart}. & Falha de gravação não deve confirmar progresso de aula. \\
\hline
ViewModels & Coordenação de estado, \path{test/unit/workspace_test.dart} e \path{test/compilation_failure_test.dart}. & Falha de compilação/ponte tratada no estado, sem sucesso fictício. \\
\hline
FFI & Biblioteca nativa real, \path{test/engine_test.dart}. & Modelo com erro esperado e posições após caracteres acentuados. \\
\hline
Widgets & Editor, navegação e trilha, arquivos de testes em \path{test/}. & Importação cancelada/com falha preserva rascunho; rota de aula bloqueada não libera etapa. \\
\hline
\end{tabular}
\vspace{0.2cm}
\fonte{O autor (2026).}
\end{quadro}
\FloatBarrier

\subsection{Verificação complementar de 15 de novembro de 2026}
\label{sec:verificacao_atual}

Em 15/11/2026, o motor foi compilado com
\texttt{cargo build --offline --manifest-path rust\_core/Cargo.toml} e testado com
\texttt{cargo test --offline --manifest-path rust\_core/Cargo.toml}.
Foram aprovados 68 testes Rust, sem falhas. A suíte Flutter foi executada com a biblioteca nativa compilada, por
\texttt{flutter test --no-pub --coverage --reporter expanded}: 159 testes foram
aprovados. \texttt{flutter analyze --no-pub} informou ``No issues found!''.

Esses números se referem ao código \texttt{0c13c2f}, não à versão do APK da
coleta. Os valores históricos de 65 Rust e 155 Flutter na
Tabela~\ref{tab:validacao_tecnica} permanecem referentes a 11/11/2026.
O motor registrava 65 testes tanto em outubro quanto em novembro; a contagem
não cresceu entre esses dois marcos. Os três testes adicionais da versão atual
tratam da verificação semântica de divisão por zero, acrescentada posteriormente.
Contagem de casos não mede, por si, extensão ou qualidade da cobertura.

O arquivo \path{coverage/lcov.info} contém 71 arquivos instrumentados, com
3.764 de 4.569 linhas executadas ao menos uma vez, correspondentes a 82,38\%.
Como recorte adicional explícito, excluindo \path{lib/src/rust/}, onde estão os
ligações geradas, restam 63 arquivos e 3.548 de 3.917 linhas, ou 90,58\%.
O primeiro é o agregado de todos os arquivos presentes no relatório LCOV (formato de relatório de cobertura de código); o segundo não
substitui esse total. São medidas de linhas Dart nos arquivos registrados,
sem cobertura de ramos, do código Rust ou de execução em Android físico.

A cobertura do motor Rust foi obtida com \texttt{cargo-llvm-cov} 0.9.1 sobre o
compilador \texttt{rustc} 1.102.0-nightly, exigido pela instrumentação de ramos, por
\texttt{cargo +nightly llvm-cov --branch --lib --tests}, executando os mesmos 68
testes. Foram excluídos o arquivo de ligações gerado (\path{frb_generated.rs}) e a
CLI (\path{src/bin/cli.rs}), que não possui testes; com a CLI, as linhas cobertas
caem para 74,38\%. No recorte do motor, foram executadas 2.053 de 2.611 linhas
(78,63\%), 207 de 225 funções (92,00\%), 2.269 de 2.862 regiões (79,28\%) e 156
de 222 ramos (70,27\%). Os menores valores de linhas estão no interpretador
(66,90\%) e na representação da AST para visualização (66,02\%); o analisador
sintático registra 80,11\% e o semântico, 82,12\%. A instrumentação de ramos é
um recurso instável do compilador \textit{nightly}. Os relatórios integrais foram gerados em texto, LCOV e HTML.

Além das suítes, foram executados 33 programas pela CLI, com registro de fonte,
saída, diagnósticos, AST e rastro retornados. Esses ensaios documentam os exemplos
deste capítulo; não foram adicionados como testes permanentes ao aplicativo.
A avaliação do artefato deve manter separadas essa verificação atual, a
verificação histórica e as percepções do levantamento.

\subsection{Protocolo e resultados das medições}
\label{sec:medicoes_pendentes}

As Tabelas~\ref{tab:medicoes_motor} e~\ref{tab:medicoes_apk} registram as
medições realizadas, com identificação do ambiente, do hash do APK, da revisão do
motor, do modo de compilação e da ferramenta de medição.

Para tempo de análise e interpretação, foram gerados programas de 10, 100 e
1.000 linhas: uma declaração \texttt{var x = 0;}, $L-2$ atribuições
\texttt{x = x + 1;} e uma impressão final. A quantidade de variáveis fica fixa.
A análise cronometra lexer, parser e semântico; a interpretação é
cronometrada separadamente sobre a AST pronta, com compilação de distribuição,
dez execuções de aquecimento e trinta amostras por tamanho, registrando mediana
e intervalo observado. Tempo total de um processo CLI não substitui essas etapas.

Para memória do rastro, foram usados programas sem laço, com uma declaração e
atribuições até totalizar 100, 1.000 e 10.000 eventos registrados, mantendo uma
variável e valores de tamanho limitado, com conferência do número de passos no
resultado. A memória atribuída ao rastro foi distinguida da memória do processo; o
tamanho do conjunto residente (RSS, do inglês \textit{Resident Set Size}) e o
\textit{heap} Dart não foram medidos.

O protocolo do motor foi executado pelo programa
\path{rust_core/examples/medicoes.rs}, com
\texttt{cargo run --release --example medicoes}, em 15/11/2026, sobre a revisão
\texttt{0c13c2f}. O ambiente foi um notebook com processador AMD Ryzen 7 4800HS,
16 GB de memória, Fedora Linux (núcleo 7.3, x86\_64) e \texttt{rustc} 1.96.0. Os
tempos foram tomados com \texttt{std::time::Instant}; a análise inclui
\texttt{tokenize}, \texttt{parse} e \texttt{analyze}, e a interpretação parte da
AST pronta. A memória foi medida por um alocador global que contabiliza bytes de
\textit{heap} alocados e liberados: ``retido'' é o acréscimo ainda vivo após a
interpretação, quando o resultado, incluindo o rastro, permanece em memória; ``pico''
é o maior acréscimo durante a execução. Não se trata de RSS do processo nem de
\textit{heap} Dart. Cada programa produziu exatamente o número de passos previsto,
sem erros; o caso de 10.000 passos coincide com o limite padrão do interpretador.

\begin{table}[htb]
\centering
\legenda[tab:medicoes_motor]{Medições de tempo e memória do motor em computador de mesa}
\footnotesize
\begin{tabular}{P{4.3cm}P{2.2cm}P{3.0cm}P{3.2cm}}
\toprule
Ensaio & Escala & Mediana ou retido & Mínimo--máximo ou pico \\
\midrule
Análise (ms) & 10 linhas & 0,013 & 0,013--0,030 \\
Análise (ms) & 100 linhas & 0,132 & 0,131--0,136 \\
Análise (ms) & 1.000 linhas & 1,334 & 1,325--1,360 \\
Interpretação (ms) & 10 linhas & 0,008 & 0,008--0,010 \\
Interpretação (ms) & 100 linhas & 0,071 & 0,071--0,165 \\
Interpretação (ms) & 1.000 linhas & 0,687 & 0,677--0,716 \\
Memória adicional (bytes) & 100 passos & 205.054 & 205.395 \\
Memória adicional (bytes) & 1.000 passos & 1.997.086 & 1.997.427 \\
Memória adicional (bytes) & 10.000 passos & 21.164.206 & 21.164.547 \\
\bottomrule
\end{tabular}
\vspace{0.2cm}
\fonte{O autor (2026). Tempos: 10 aquecimentos e 30 amostras por tamanho.
Memória: uma execução por tamanho.}
\end{table}
\FloatBarrier

Nessas condições, análise e interpretação cresceram de forma aproximadamente
linear com o número de linhas, e o rastro reteve cerca de 2 kB por passo com uma
única variável, pois cada passo guarda um retrato das variáveis e dos escopos. Os
valores descrevem o motor em um computador de mesa e não estimam o desempenho em
celulares, que não foi medido.

O tamanho do APK foi obtido em bytes do arquivo identificado por SHA-256, e a
versão mínima foi lida do manifesto empacotado, sem substituí-la pelo valor
presumido da configuração Gradle.

As medidas do pacote referem-se ao APK publicado em
\url{https://luminthor.vercel.app/}, obtido em 15/11/2026 e lido com
\texttt{aapt dump badging} (Android \textit{build-tools} 36.1.0). O pacote contém bibliotecas nativas para \texttt{arm64-v8a},
\texttt{armeabi-v7a} e \texttt{x86\_64}, o que contribui para seu tamanho.

\begin{table}[htb]
\centering
\legenda[tab:medicoes_apk]{Identificação e medições do pacote Android}
\footnotesize
\begin{tabular}{P{3.6cm}P{3.0cm}P{6.1cm}}
\toprule
Métrica & Unidade/registro & Resultado \\
\midrule
Tamanho do APK & Bytes & 63.692.511 \\
Android mínimo do APK & Nível de API e versão & API 24 (Android 7.0); alvo API 36 \\
Identificação do pacote & Versão e SHA-256 & \texttt{com.ide.ide}, versão 1.0.0 (código 1); \texttt{fd5cae4a239ca0b3\-abae886a0351d3af\-74fb4f6bb825339f\-e2599aec62e9174a} \\
\bottomrule
\end{tabular}
\vspace{0.2cm}
\fonte{O autor (2026).}
\end{table}
\FloatBarrier

\subsection{Identificação do código-fonte}
\label{sec:repositorio_luminthor}

O código-fonte está disponível no repositório público
\url{https://github.com/Vinni071/Snuppy}. O marco de 11/11/2026 corresponde ao
commit \texttt{d2fbb0258ea2b4b21d95baac6796a1d8a83590ac}, e a versão descrita
neste trabalho, ao commit \texttt{0c13c2f}.
